%----------------------------------------------------------------------
\section{Introduction}
%----------------------------------------------------------------------

Zermelo--Fraenkel set theory with Choice (ZFC) takes membership
($\in$) as an undefined primitive and regulates it by axioms.
The axiom of separation, for instance, exists primarily to prevent
the formation of Russell's set. The law of excluded middle is
inherited from the ambient classical logic. Boolean complementation
--- $x \notin A$ if and only if $x \in A^c$ --- is available by default.

We propose an alternative starting point: \emph{discrimination}.
Rather than asking whether $x \in A$ as a pre-given fact, we ask
whether a discriminator $d$ can separate $x$ from the rest of a
collection $S$ under a discrimination regime $\kappa$. The key
departure from the classical account is:

\begin{enumerate}[label=(\roman*)]
  \item Membership is a three-place relation: element, collection,
    discriminator.
  \item The algebra of discriminators is the exterior algebra
    $\bigwedge(\GF(2)^n)$, not Boolean algebra.
  \item Boolean complementation is not given; it is an achievement,
    local to a dimension of $\kappa$ where exhaustiveness and
    exclusivity have been established.
  \item The discrimination structure $\kappa$ grows monotonically:
    articulating a new predicate \emph{is} a $\kappa$-evolution.
    The observer is not independent of the system.
\end{enumerate}

The framework is guided by the following principle, which
we adopt as an explicit foundational axiom --- comparable
in status to ZFC's axiom of infinity or the law of excluded
middle in classical logic:

\begin{definition}[Operationalist axiom]\label{def:operationalist}
If a distinction cannot be witnessed by any discriminator
in any $\kappa$-state, it is not a valid distinction within
the framework. Equivalently: if no $\kappa$-evolution
produces a discriminator that separates two claims, the
claims are $\kappa$-equivalent and hence identical by
extensionality (Theorem \ref{thm:ext}).
\end{definition}

In informal terms: \emph{if a distinction is real, it must
be constructible; if constructible, it must be behaviorally
reachable; if reachable, it must be observable; if
observable, it must be coverable; if not, it is not a valid
runtime distinction.} This is a \emph{choice} about what
the framework admits as meaningful, with strong engineering
motivation (a framework that admits unobservable distinctions
cannot be implemented, tested, or falsified), but it is not
a forced mathematical conclusion.

\subsection{Reading guide: how to verify this paper's claims}
\label{sec:reading-guide}

The paper contains 125 formal objects across 10 sections.
Before encountering them, the reader should know three
things: where each claim is verified, what is not claimed,
and how much of the paper depends on which axioms.

\medskip
\noindent\textbf{Axiom-class distribution.}
The paper rests on three commitments beyond the algebra:
profile linearity (\textsf{P}), evaluation linearity
(\textsf{E}), and the operationalist axiom (\textsf{O}).
Of the 125 formal objects: 101 (81\%) depend only on
the algebra. 9 additionally require profile linearity.
6 require both linearities. 8 require the operationalist
axiom. 1 requires all three. The operationalist axiom
--- the most philosophically contentious commitment ---
touches 6\% of the paper's formal content.

\medskip
\noindent\textbf{Verification map.}
Each theorem's proof lives in the spine (\S\S3--9); each
theorem's component-by-component check lives in Appendix
B. Toulmin footnotes at each theorem state: axiom class,
dependency depth, verification pointer, warrant, qualifier,
finite-$\kappa$ validity, and novelty status.

\begin{center}
\scriptsize
\begin{tabular}{l|c|c|l}
\textbf{Claim} & \textbf{Class} & \textbf{Status} &
  \textbf{Verification} \\
\hline
Extensionality (4.2) & \textsf{A} & \cmark &
  Tautological by construction \\
Russell exclusion (4.5) & \textsf{AEP} & \cmark &
  B.2; affine offset explicit \\
Foundation (4.13) & \textsf{AP} & \cmark &
  B.16, B.34; depth $\leq \dim(V)$ \\
Choice, finite (4.33i) & \textsf{A} & \cmark &
  B.17, B.31; measure $= \dim(V/\kappa)$ \\
Choice, infinite (4.33ii) & \textsf{AO} & \omark &
  Axiom-dependent [O3] \\
Indistinguishability (4.26) & \textsf{AO} & \cmark &
  B.35; stabilization + linearity \\
Exhaustivity (5.7) & \textsf{A} & \cmark &
  B.18; filling = $\kappa$-evolution (Leibniz) \\
Chain complex (5.5) & \textsf{A} & \cmark &
  Double-sum cancellation explicit \\
Groupoid (5.10) & \textsf{A} & \cmark &
  B.19; each axiom checked \\
Univalence (5.12) & \textsf{AO} & \cmark &
  B.36; $\bigcap \kappa^\perp = \{0\}$ \\
2-cat axioms (6.8) & \textsf{A} & \cmark &
  B.20; interchange via GF(2)$^2$ profiles \\
Adjunction (6.11) & \textsf{A} & \cmark &
  B.1; all components verified \\
Yoneda (6.17) & \textsf{A} & \cmark &
  B.21; G7 resolved (presheaf boundary) \\
Limits (6.20) & \textsf{A} & \cmark &
  B.26; all five types enumerated \\
$(n,1)$-structure (7.1) & \textsf{A} & \cmark &
  B.29; commutativity of $\wedge$ $\Rightarrow$ self-inverse \\
Dynamics transfer (7.12) & \textsf{A} & \cmark &
  B.30; linearities swap under $S_3$ \\
Free $(n,k)$ (7.15) & \textsf{A} & \cmark &
  B.37; groupoid + Grassmannian \\
Sym.\ monoidal (8.2) & \textsf{A} & \cmark &
  Strict: $\wedge$ literally assoc., all coherences id \\
Delooping (8.4) & \textsf{A} & \cmark &
  B.22; rotation explicit \\
Periodic table (8.8) & \textsf{A} & \cmark &
  B.7, B.39; recognizer rotation \\
Mon.\ closure (8.13) & \textsf{A} & \cmark &
  B.22; naturality via strict unit \\
Self-enrichment (9.2) & \textsf{A} & \cmark &
  B.8; all components checked \\
$\Omega$ classifier (9.4) & \textsf{A} & \cmark &
  B.9, B.27; composites via linearity \\
Fibered topos (9.6) & \textsf{AP} & \cmark &
  B.4, B.14, B.27, B.38; O4 resolved \\
Proof theory (9.11) & \textsf{A} & \cmark &
  B.10; truth tables explicit \\
FOL internal (9.12) & \textsf{AP} & \cmark &
  B.15; all rules constructed \\
Self-hosting (9.14) & \textsf{AP} & \cmark &
  B.11; tables explicit \\
Fixed point (9.15) & \textsf{A} & \cmark &
  B.12, B.33; unique (1-dim), $S_3$-inv \\
Irremovable (9.16) & \textsf{AP} & \cmark &
  B.13; O5 resolved (dim = 1) \\
Sheaf topos (9.18) & \textsf{AP} & \cmark &
  B.14, B.28; topology enumerated \\
Fano closure (9.20) & \textsf{A} & \cmark &
  B.5; all 7 lines explicit \\
Convergence (9.22) & \textsf{A} & \cmark &
  B.23; standard finite geometry \\
Pairing (4.18) & \textsf{A} & \cmark &
  B.32; equiv.\ class, not char.\ function
\end{tabular}
\end{center}

\noindent Legend: \cmark\ = verified,
\gmark\ = gap (specific issue identified),
\omark\ = open-acknowledged.

\medskip
\noindent\textbf{What is not claimed (open questions).}

\emph{Open questions} (acknowledged, not claimed):
\begin{description}[nosep,leftmargin=0.8cm]
  \item[O1--O6] All resolved except O3 (by design):
    $\in$-induction, infinite choice, join-semilattice,
    fixed-point uniqueness [\textbf{resolved}: $\dim = 1$],
    join-semilattice [\textbf{resolved}: $V^*$ is a
    vector space $\Rightarrow$ subspace lattice], CD hexagon axioms.
\end{description}

\medskip
\noindent\textbf{Gap priorities} (from SPPF downstream
impact analysis).

\emph{Tier 1 --- load-bearing} (closing these unblocks the
most downstream results):
\begin{description}[nosep,leftmargin=0.8cm]
  \item[G6] Interchange law (10 downstream).
    \textbf{Closed.} Cross terms cancel via $1+1=0$;
    grade-4 collapses to grade 2 in $\GF(2)^2$.
    Unblocks G8, G16.
  \item[G4] Exhaustivity filling (8 downstream).
    \textbf{Closed.} Leibniz rule: $\partial(R \wedge
    e_{n+1}) = R$. Unblocks G2, G8, G16.
  \item[G12] Symmetric monoidal (8 downstream).
    \textbf{Closed.} Strict monoidal: all coherences
    are identities. Unblocks G13, G16.
\end{description}

\emph{Tier 2 --- structural} (unblocks \S9 claims):
\begin{description}[nosep,leftmargin=0.8cm]
  \item[G13] Monoidal closure naturality (6 downstream).
    \textbf{Closed.} Strict unit: $d_h \wedge 1 = d_h$.
    Unblocks self-enrichment, topos, FOL.
  \item[G8] Limits (5 downstream). \textbf{Closed.}
    All five limit types enumerated; universality from
    $\GF(2)$-linearity. Unblocks topos.
  \item[G14] Power objects (5 downstream).
    \textbf{Closed.} Six cases; composites by linearity;
    Galois-equivariant transitions. Unblocks topos
    (O4 resolved).
\end{description}

\emph{Tier 3 --- localized} (1--2 downstream results):
G3 (indistinguishability $\omega$-chain), G17 (pairing
linearity), G1 (foundation chain bound), G2 (choice
termination), G9/G10 (higher structure), G15 (fixed-point
stability).

\emph{Tier 4 --- terminal or downgradable} (0 downstream):
G5 (univalence: lift from analogy to native
univalence), G11 (free $(n,k)$ coherence),
G16 (sheaf topology non-triviality).

\emph{All 17 gaps closed. All 33 audited claims verified
(see Appendix D for the technique matrix and rotation principle).}

\emph{Identified gaps} (claimed but incompletely verified):
\begin{description}[nosep,leftmargin=0.8cm]
  \item[G1] Foundation chain bound: \textbf{closed}.
    Each link $a_i \ni a_{i+1}$ produces a link vector
    $v_i = a_i + a_{i+1} \in \GF(2)^n$. A dependent
    link vector is a derived relation, not a new
    membership step. Chain depth $=
    \dim(\mathrm{span}\{v_i\}) \leq n = \dim(V)$.
    Sharp: achieved by basis chain. Galois: each
    independent link reduces $\dim(V/\kappa)$ by 1.
    (\S4, Thm 4.13; Appendix B.34.)
  \item[G2] Finite choice termination: \textbf{closed}.
    Decreasing measure $m(\kappa) = \dim(V/\kappa)$.
    Each evolution adds one discriminator:
    $m(\kappa') = m(\kappa) - 1$. Terminates in at most
    $n = \dim(V)$ steps. Evolutions are global (shared
    across sub-families). At $m = 0$, Galois group is
    trivial, all classes are singletons, selection trivial.
    (\S4, Thm 4.33; Appendix B.31.)
  \item[G3] Indistinguishability $\omega$-chain:
    \textbf{closed}. Three steps: (1) ascending chain
    stabilizes ($V$ finite-dim), (2) limit regime is
    finite span of individual $d_i$'s, (3) profile
    linearity: linear combination of non-separating
    discriminators is non-separating. Galois: residual
    symmetry group contains the unbreakable symmetries.
    (\S4, Prop 4.26; Appendix B.35.)
  \item[G4] Exhaustivity: \textbf{closed}. The filling
    is $C = R \wedge e_{n+1}$ for any independent $e_{n+1}$.
    Leibniz: $\partial(R \wedge e_{n+1}) = \partial R
    \wedge e_{n+1} + R \cdot 1 = R$. $\kappa$-evolution
    itself is the constructive filling procedure.
    (\S5, Thm 5.7; Appendix B.18.)
  \item[G5] Native univalence: \textbf{closed}. The
    path space is $\kappa^\perp$ (the annihilator). The
    operationalist axiom says $\bigcup \kappa = V$,
    so $\bigcap \kappa^\perp = V^\perp = \{0\}$.
    Persistent equivalence ($a + b \in \bigcap
    \kappa^\perp$) implies $a + b = 0$, i.e.\ $a = b$.
    The ua map and idtoeqv are inverses. Path space is
    contractible (single point). Univalence is
    Galois-group exhaustion.
    (\S5, Prop 5.12; Appendix B.36.)
  \item[G6] Interchange law: \textbf{closed}.
    Both paths through the interchange square yield
    $p(g \circ f) + p(g' \circ f') \in \GF(2)^2$;
    cross terms cancel because $1+1=0$. The grade-4
    expression collapses to grade 2 via the Klein
    four-group's characteristic 2. (\S6, Thm 6.8;
    Appendix B.20.)
  \item[G7] Yoneda, infinite case: \textbf{resolved}
    via Galois structure. At each finite $\kappa$, Yoneda
    holds (verified). The pro-system of finite Yoneda
    lemmas has transition maps (restriction functors
    between fibers) that are equivariant under the Galois
    group inclusion $\mathrm{Gal}(V/\kappa) \hookrightarrow
    \mathrm{Gal}(V/\kappa')$ for $\kappa' \subseteq \kappa$.
    Unparsed data is the Galois group's orbit: data whose
    categorical structure is undetermined by the current
    $\kappa$. The orbit shrinks as $\kappa$ grows (Galois
    group reduction). ``Infinite Yoneda'' is the inverse
    limit of this Galois-equivariant pro-system, not a
    single lemma applied to an infinite category.
    (\S6, Thm 6.17; \S\ref{sec:galois}.)
  \item[G8] Limits: \textbf{closed}. Exhaustive
    construction of terminal object, product, equalizer,
    pullback, and general finite limit. Universality
    from linearity over $\GF(2)$: factoring morphisms
    are unique because linear maps are determined by
    basis values. Filling via $\kappa$-evolution when
    constraints require new discriminators.
    (\S6, Prop 6.20; Appendix B.26.)
  \item[G9] $(n,1)$-structure: \textbf{closed}. Over
    $\GF(2)$, $\wedge$ is commutative, so every grade-$k$
    element ($k \geq 2$) equals its own reversal:
    $d_k \wedge \cdots \wedge d_1 = d_1 \wedge \cdots
    \wedge d_k$. Every $k$-morphism is self-inverse.
    Direction requires asymmetry, which only the
    $\tau/\sigma$ population interface provides, and
    that interface exists only at grade 1.
    (\S7, Thm 7.1; Appendix B.29.)
  \item[G10] Dynamics transfer: \textbf{closed}. The
    evaluation map $\mathrm{ev}(d,a) = d(a)$ is bilinear.
    $S_3$ rotation swaps arguments:
    $\mathrm{ev}'(a,d) = \mathrm{ev}(d,a)$. Profile
    linearity and evaluation linearity \emph{swap} under
    rotation; since both hold, both are preserved.
    The bilinear pairing is $GL$-invariant; $S_3 \subset GL$.
    (\S7, Prop 7.12; Appendix B.30.)
  \item[G11] Free $(n,k)$ Segal conditions:
    \textbf{closed}. Groupoids automatically satisfy
    Segal conditions (unique fillers). The $k$-cells
    are the Grassmannian $\mathrm{Gr}(k, V)$ over
    $\GF(2)$: decomposable $k$-forms, each determined
    by $k$ independent 1-forms. The Segal map is a
    bijection. Non-decomposable elements are chain-group
    elements, not $k$-morphisms.
    (\S7, Thm 7.15; Appendix B.37.)
  \item[G12] Symmetric monoidal: \textbf{closed}.
    $\wedge$ is strictly associative with strict unit;
    all coherence morphisms (associator, unitors,
    braiding) are identities. Pentagon, triangle, hexagon
    trivially satisfied. Galois mirror: joins of
    $\kappa$-states are strict; Galois group reduction
    under joins = intersection of subgroups.
    (\S8, Prop 8.2; Appendix B.22.)
  \item[G13] Monoidal closure naturality: \textbf{closed}.
    Precomposition with $h \otimes \mathrm{id}_B =
    d_h \wedge 1 = d_h$ (strict unit). Naturality square
    commutes because restricting $d_h \wedge 1$ to the
    $A$-component yields $d_h$. Galois: functoriality of
    restriction along fiber inclusions.
    (\S8, Thm 8.13; Appendix B.22.)
  \item[G14] Power objects: \textbf{closed}. Six cases
    enumerated: single-discriminator, linear combination,
    wedge product, general at stage $\kappa$, at joins,
    and the Galois perspective (equivariant transition
    functors under Galois reduction). Subobjects form a
    linear subspace of $\operatorname{Hom}(B, \Omega)$;
    uniqueness from $\GF(2)$-linearity. The fibered
    colimit classifies every subobject classifiable at
    any stage.
    (\S9, Thms 9.4, 9.6; Appendix B.27.)
  \item[G15] Fixed-point stability: \textbf{closed}.
    $\bigwedge^2(\GF(2)^2)$ is 1-dimensional: exactly
    one non-zero 2-cell exists ($e_1 \wedge e_2$).
    No alternative to perturb to. Perturbation to zero
    triggers irremovability (reinstatement). The 2-cell
    is $GL(2,\GF(2))$-invariant ($\det = 1$ over
    $\GF(2)$). Also resolves [O5]: uniqueness from
    dimensionality, not $H_2 = 0$.
    (\S9, Thm 9.15; Appendix B.33.)
  \item[G16] Sheaf topology: \textbf{closed}. Exhaustive
    enumeration: single-axis families (3, each 1-dim)
    do not cover. Two-axis families (12, each spanning
    $\GF(2)^2$) cover. Not trivial (proper subfamilies
    cover), not discrete (single-axis sieves fail).
    The topology is the unique $S_3$-invariant topology
    determined by the triangle identity.
    (\S9, Thm 9.18; Appendix B.28.)
  \item[G17] Pairing: \textbf{closed}. The pair
    $\{a,b\}$ is a $\kappa$-equivalence class, not a
    characteristic function. Condition: $a + b \in
    \kappa^\perp$ (annihilator). No linear ``is $a$
    or is $b$'' predicate needed. The pair is defined
    by what $\kappa$ \emph{cannot} distinguish.
    Verified in concrete model.
    (\S4, Prop 4.18; Appendix B.32.)
\end{description}

\medskip
\noindent\textbf{Gap priorities} (from SPPF downstream
impact analysis).



